\SourceLeaf{336}

\section*{[II.]}
\begin{center}PREPARATIONS FOR THE GENERAL THEORY\\OF BIQUADRATIC RESIDUES.\end{center}
\pdfbookmark[1]{II. Preparations for the General Theory of Biquadratic Residues}{nachlass-biquadratic-II}
\BookRule{.15\textwidth}
\par\medskip\pdfbookmark[2]{(1.)}{nachlass-biquadratic-II-prefatory-1}\centerline{(1.)}\smallskip
Let $P=x+iy$, where neither $x$ nor $y$ is an integer. We denote the number $+1$ by $LP,L'P,L''P,L'''P$, according as $P$ lies in the first, second, third or fourth quadrant (in the first and second quadrants $[y]$ is even, in the third and fourth odd; in the first and fourth $[x]$ is even, in the second and third odd). In all cases where these symbols are not $=1$, they are assumed to be $=0$. We then have the following $24$ relations:

\[
\begin{array}{lll}
L(P\pm1)=L'P&L(P\pm i)=L'''P&L(P\pm1\pm i)=L''P\\
L'(P\pm1)=LP&L'(P\pm i)=L''P&L'(P\pm1\pm i)=L'''P\\
L''(P\pm1)=L'''P&L''(P\pm i)=L'P&L''(P\pm1\pm i)=LP\\
L'''(P\pm1)=L''P&L'''(P\pm i)=LP&L'''(P\pm1\pm i)=L'P
\end{array}
\]
\[
\begin{array}{lll}
LiP=L'''P&L(-P)=L''P&L(-iP)=L'P\\
L'iP=LP&L'(-P)=L'''P&L'(-iP)=L''P\\
L''iP=L'P&L''(-P)=LP&L''(-iP)=L'''P\\
L'''iP=L''P&L'''(-P)=L'P&L'''(-iP)=LP
\end{array}
\]

\par\medskip\pdfbookmark[2]{(2.)}{nachlass-biquadratic-II-prefatory-2}\centerline{(2.)}\smallskip
By $PP'$ or $z$ we denote a line beginning at $P$ and ending at $P'$. It need not be straight. We assign a positive and a negative side to every straight line drawn from $2x+2iy$ to $2x+(2y+1)i$ (where $x,y$ range over all integers); we choose the right side as positive and the left as negative. By $Tz$ we denote the number of all intersections
\SourceLeaf{337}

of the line $z$ with the lines just described, counting positively those where $z$ passes from the negative side to the positive, and negatively the others. We also put
\[
Tz-T(z-1)=Sz
\]

\Needspace{18\baselineskip}
($z-1$ is a line parallel to $z$, going from the point $P-1$ to $P'-1$). Evidently we need only add to the above system the lines drawn from $2x+1+2yi$ to $2x+1+(2y+1)i$, regarding their left sides as positive and their right sides as negative, to recognize in $Sz$ the number of all intersections of $z$ with this double system of straight lines. We further have
\[
\begin{aligned}
T(-z)&=-T(z+i)\\
T(z)+T(z+i)&=[\tfrac12x]-[\tfrac12x']\\
S(z+1)&=-Sz\\
S(z+i)&=-Sz+LP+L'''P-LP'-L'''P'\\
S(z+1+i)&=Sz-LP+L'''P+LP'+L'''P'\\
Siz&=Sz-LP+LP'\\
S(-z)&=Sz-LP-L'''P+LP'+L'''P'\\
S(-iz)&=Sz+L'P-L'P'
\end{aligned}
\]

\renewcommand{\Article}[1]{\par\medskip\Needspace{4\baselineskip}\pdfbookmark[2]{Art. #1}{nachlass-biquadratic-II-art-#1}\centerline{\textbf{#1.}}\nobreak\smallskip}
\Article{1}
We consider two kinds of points in the plane: first, those corresponding to integers; second, those determined by multiplying integers by the quantity $Q=\frac{m}{2M}$. We can distinguish them by calling them points of the first order and points of the second order.

\Article{2}
Joining each point of the second order to its four neighbours by straight lines, which we shall call \emph{ligatures}, divides the whole plane into infinitely many squares. The points of the first order lie partly inside these squares, partly on the ligatures between their
\SourceLeaf{338}

endpoints, and partly at the endpoints of the ligatures, the last case occurring when they are also points of the second order. If $kQ$ is such a point, then, provided $m,M$ have no common divisor and are both odd, $k$ must be divisible by $M$.

\Article{3}
For ligatures we can also distinguish the starting point from the endpoint, and thus $PQ$ from $QP$, or in some cases set this distinction aside. We call two such ligatures opposite. In general, the most convenient way to denote ligatures is by their starting points and endpoints, which may if desired be enclosed in brackets. The ligature opposite to a given one shall be indicated by a double overline: $QP=\overline{\overline{PQ}}$.

\Needspace{13\baselineskip}
\Article{4}
Each of the squares in question is enclosed by four such ligatures
\[
\begin{gathered}
\{kQ,(k+1)Q\},\ \{(k+1)Q,(k+1+i)Q\},\\
\{(k+1+i)Q,(k+i)Q\},\ \{(k+i)Q,kQ\}\dotsfour\Omega
\end{gathered}
\]
to whose right it lies. Here it is important to consider the form of the number $k$. We distinguish four kinds of squares according as $k\equiv0,1,1+i,i\ (\mathrm{mod.}\ 2)$, and use the numbers $0,1,2,3$, which we call the respective intensors of the squares.

\Article{5}
We assign to the ligatures the same intensors as those of the squares lying to their right.

\Article{6}
We must now consider another, larger square $\Omega'$, namely the one obtained when the square indicated in $4$ for $k=0$ is multiplied by $M$: it is therefore bounded by the straight lines $\mu,\mu',\mu'',\mu'''$
\[
\{0,\tfrac12m\},\ \{\tfrac12m,\tfrac12(1+i)m\},\ \{\tfrac12(1+i)m,\tfrac12im\},\ \{\tfrac12im,0\}
\]
It consists of complete squares $\Omega$ and pieces of such squares. Add up
\SourceLeaf{339}

all points of the first order inside it, taking for each point the intensor of the square $\Omega$ in which it lies. This sum, or its least residue modulo $4$, is called the decident of $M$ for the modulus $m$, and determines the biquadratic modality of $M$ with respect to this modulus.

\Article{7}
We divide the square $\Omega'$ into $5$ pieces as follows. Join the point $0$ to $\frac14(1+i)(m-1)$ by the line $\lambda$, which is to run entirely along ligatures inside $\Omega'$. Let
\[
\tfrac12m+i\lambda=\lambda',\quad\tfrac12(1+i)m-\lambda=\lambda'',\quad\tfrac12im-i\lambda=\lambda'''
\]
These $4$ lines therefore run from the corners of the square $\Omega'$ into its interior and end at the four corners of the innermost square. Its intensor will be $0$ if $m\equiv1\ (\mathrm{mod.}\ 2+2i)$; let the ligatures of this square be $\nu,\nu',\nu'',\nu'''$.

The $5$ pieces will therefore be bounded by
\[
\begin{array}{rll}
\mathrm{I}.&\ldots&\mu,\ \lambda',\ \overline{\overline\nu},\ \overline{\overline\lambda}\\
\mathrm{II}.&\ldots&\mu',\ \lambda'',\ \overline{\overline{\nu'}},\ \overline{\overline{\lambda'}}\\
\mathrm{III}.&\ldots&\mu'',\ \lambda''',\ \overline{\overline{\nu''}},\ \overline{\overline{\lambda''}}\\
\mathrm{IV}.&\ldots&\mu''',\ \lambda,\ \overline{\overline{\nu'''}},\ \overline{\overline{\lambda'''}}\\
\mathrm{V}.&\multicolumn{2}{l}{\text{the inner square }\nu,\ \nu',\ \nu'',\ \nu'''}
\end{array}
\]
The decident is therefore the enumeration of all points of the first order in I. II. III. IV.

\Article{8}
For brevity, the intensor of any point shall mean the intensor of the square in which it lies, and shall be expressed by a prefixed $\Upsilon$.

\Article{9}
The decident is therefore
\[
\Sigma\Upsilon P+\Sigma\Upsilon P'+\Sigma\Upsilon P''+\Sigma\Upsilon P'''
\]
where $P$ denotes all points in I., and so on.

\SourceLeaf{340}

\Article{10}
We now also consider the region $\mathrm{VI}=-i\mathrm{IV}$, which lies outside $\Omega'$, but adjoins I along $\mu$ and together with it forms the region $\omega$, consisting of $AA+BB$ complete squares. If $\Pi$ denotes all integer points, and $\Pi'$ all integer points increased by $\frac12i$, in this region, it is easy to prove that the decident
\[
\begin{aligned}
={}&\Sigma\Upsilon\Pi-\Sigma\Upsilon\Pi'+\text{number of all integer points inside VI}\\
&-\text{number of all half-points inside VI.}
\end{aligned}
\]
